% Context from chapters/machine-learning.tex; for mathematical reference, not compilation.
c{\la}{2} \norm{\be}_\Hh^2
}

Start with regression under the squared loss $\ell(y,z)=\frac{1}{2}(y-z)^2$:
\eq{
	\be^\star = \uargmin{\be \in \Hh} \norm{\Phi\be-y}_{\RR^n}^2+\la\norm{\be}_\Hh^2.
}
which is the ridge regression problem studied in Section~\ref{sec-linear-models}. 
 
Setting the gradient to zero gives
\eq{
	\be^\star = ( \Phi^*\Phi + \la \Id_{\Hh} )^{-1} \Phi^* y.
}
Direct computation may be impractical or impossible when $\Hh$ is infinite dimensional. The following Woodbury identity converts the problem into a finite system.
 
This is the feature-space counterpart of~\eqref{eq-linest-woodbury}, with $\Phi$ replacing $X$ and $\la$ replacing $n\la_n$.

\begin{prop}
	One has
	\eq{
		( \Phi^*\Phi + \la \Id_{\Hh} )^{-1} \Phi^* = \Phi^* ( \Phi \Phi^* + \la \Id_{\RR^n} )^{-1}. 
	}
\end{prop}
\begin{proof}
	Denoting $U=( \Phi^*\Phi + \la \Id_{\Hh} )^{-1} \Phi^*$ and $V=\Phi^* ( \Phi \Phi^* + \la \Id_{\RR^n} )^{-1}$, one has
 	\eq{
		( \Phi^*\Phi + \la \Id_{\Hh} ) U = \Phi^*
	}
	and
	\eq{
		( \Phi^*\Phi + \la \Id_{\Hh} ) V = ( \Phi^*\Phi\Phi^* + \la \Phi^*)  ( \Phi \Phi^* + \la \Id_{\RR^n} )^{-1}
		 = \Phi^* ( \Phi \Phi^* + \la \Id_{\RR^n} ) ( \Phi \Phi^* + \la \Id_{\RR^n} )^{-1} = \Phi^*.		
	}
	Since $( \Phi^*\Phi + \la \Id_{\Hh} )$ is invertible, one obtains the result. 
\end{proof}

The identity gives the alternative representation
\eq{
	\be^\star = \Phi^* c^\star
	\qwhereq
		c^\star \eqdef  ( K + \la \Id_{\RR^n} )^{-1} y
}	
where the empirical kernel matrix is
\eq{
	K=\Phi\Phi^* = ( \kappa(x_i,x_j) )_{i,j=1}^n \in \RR^{n \times n}
	\qwhereq
	\kappa(x,x') \eqdef \dotp{\phi(x)}{\phi(x')}. 
}	
The coefficients $c^\star \in \RR^n$ depend on the inputs only through $\kappa$, without explicit computation of $\phi(x)$. If $\kappa$ is inexpensive to evaluate, forming $K$ takes $O(n^2)$ kernel evaluations. A dense direct solve costs $O(n^3)$ operations; an iterative solve costs $O(n^2)$ per matrix-vector product, with the iteration count determined by conditioning and accuracy.

The fitted predictor at a new input $x \in \Xx$ also uses only kernel evaluations:
\eql{\label{eq-prediction-evaluation}
	\dotp{\be^\star}{\phi(x)}_\Hh = \dotp{ \sum_i c_i^\star \phi(x_i)}{\phi(x)}_\Hh = \sum_i c_i^\star \kappa(x,x_i).
}

                                                  
\subsection{Kernel Design}

One can also start directly from a kernel $\kappa(x,x')$ rather than an explicit feature map $\phi$. This is often convenient because kernels can be designed through their geometric properties and combined, for example by addition. The required condition is that every Gram matrix $(\kappa(x_i,x_j))_{i,j}$ be symmetric positive semidefinite. Such kernels are commonly called positive definite kernels; the condition is equivalent to the existence of a Hilbert-space feature map generating the kernel.
 
The underlying input space need not be Euclidean.


When using the linear kernel $\kappa(x,y)=\dotp{x}{y}$, one retrieves the linear models studied in the previous section, and the lifting is trivial $\phi(x)=x$.
 
A family of popular kernels are polynomial ones,  $\kappa(x,x') = (\dotp{x}{x\prime}+c)^a$ for $a
\in \NN^*$ and $c>0$, which corresponds to a lifting in finite dimension. For instance, for $a=2$ and $p=2$, one has a lifting in dimension $6$
\eq{
	\kappa(x,x') = \pa{ x_1x_1'+x_2x_2' + c}^2 = \dotp{\phi(x)}{\phi(x')}
	\qwhereq
	\phi(x)=(x_1^2,x_2^2,\sqrt{2}x_1x_2,\sqrt{2c}x_1,\sqrt{2c}x_2,c)^* \in \RR^6.	
}
A widely used choice on Euclidean spaces is the Gaussian kernel
\eql{\label{eq-gauss-kernel}
	 \kappa(x,y) \eqdef e^{-\frac{\norm{x-y}^2}{2\sigma^2}} . 
}
The bandwidth $\si>0$ controls the locality of the model and is typically tuned by cross-validation. The Gaussian kernel has the infinite-dimensional feature map $\phi(x)(t)=(2/(\pi\sigma^2))^{p/4}\exp(-\norm{x-t}^2/\sigma^2)\in L^2(\RR^p)$.
 
Another related popular kernel is the Laplacian kernel $\exp(-\norm{x-y}/\si)$.
 
More generally, Bochner's theorem states that a continuous translation-invariant kernel $\kappa(x,x')=k(x-x')$ on $\RR^p$ is positive definite if and only if $k$ is the Fourier transform of a finite nonnegative measure. If this measure has an integrable density, a feature map can be built from the square root of that density. When $k$ is integrable and its Fourier transform is an integrable function, the condition is $\hat k\geq0$; strict positivity is not necessary.
 
Figure~\ref{fig-kernel} shows an example of regression using a Gaussian kernel. 


\begin{figure}
\centering
\begin{tabular}{@{}c@{\hspace{1mm}}c@{\hspace{1mm}}c@{\hspace{1mm}}c@{\hspace{1mm}}c@{}}
\includegraphics[width=.19\linewidth]{ml/regression/scatter-2d}&
\includegraphics[width=.19\linewidth]{ml/regression/kernel-2}&
\includegraphics[width=.19\linewidth]{ml/regression/kernel-3}&
\includegraphics[width=.19\linewidth]{ml/regression/kernel-4}&
\includegraphics[width=.19\linewidth]{ml/regression/kernel-5}\\
& $\si=0.1$ & $\si=0.5$ & $\si=1$ & $\si=5$
\end{tabular}
\caption{\label{fig-kernel}
Regression using a Gaussian kernel.
}
\end{figure}

\paragraph{Kernel on non-Euclidean spaces.}

On a general input space $\Xx$, a real kernel $\kappa$ must produce a symmetric positive semidefinite Gram matrix $K=(\kappa(x_i,x_j))_{i,j}$ for every finite collection of inputs. Equivalently, there is a Hilbert space $\Hh$ and a feature map $\phi:\Xx \rightarrow \Hh$ such that $\kappa(x,x')=\dotp{\phi(x)}{\phi(x')}_\Hh$.

For instance, if $\Xx = \Pp(\Mm)$ is the set of measurable sets of finite measure in a measure space $(\Mm,\mu)$, one can define $\kappa(A,B) = \mu(A \cap B)$ which corresponds to the lifting $\phi(A)=1_A$ so that $\dotp{\phi(A)}{\phi(B)} = \int 1_A(x) 1_B(x)\d \mu(x) = \mu(A \cap B)$.
 
Kernels can also be defined on strings and graphs, although their construction is more involved.

When $\kappa(x,x)>0$ for every input under consideration, the kernel can be normalized to have a unit diagonal. For a translation-invariant kernel on $\RR^p$, the diagonal equals the constant $k(0)$, so normalization only requires division by $k(0)>0$. This corresponds to replacing $\phi(x)$ by $\phi(x)/\norm{\phi(x)}_\Hh$ and thus replacing $\kappa(x,x')$ by $\frac{\kappa(x,x')}{\sqrt{\kappa(x,x)}\sqrt{\kappa(x',x')}}$.



                                                  
\subsection{General Case}

The least-squares result extends to other losses under the assumptions below. With a squared Hilbert-space norm penalty, the minimizer belongs to a subspace generated by the data, of dimension at most $n$. This representer theorem reduces the problem to finitely many coefficients even when the feature space $\Hh$ is infinite dimensional.

\begin{prop}
	Assume $\lambda>0$ and the losses $\ell(y_i,\cdot)$ are finite, nonnegative, and convex. The solution $\be^\star\in\Hh$ of
	\eql{\label{eq-kernel-generic}
		 \umin{\be \in \Hh} \sum_{i=1}^n \ell(y_i, \dotp{\phi(x_i)}{\be}) + \frac{\la}{2} \norm{\be}_\Hh^2
		 =
		 \Loss(\Phi \be,y) + \frac{\la}{2} \norm{\be}_\Hh^2
	}
	is unique and can be written as
	\eql{\label{eq-rkhs-repres